\section{Related Work}
\label{sec:related}

\paragraph{Observed-topology link prediction.}
Graph-native link predictors condition on paths, enclosing subgraphs, distance
labels, or neighborhood overlap. SEAL extracts a query-centered subgraph
\citep{zhang2018seal}; distance encoding and the labeling trick explain why
target-conditioned representations can distinguish pairs that independently
computed node embeddings cannot \citep{li2020distance,zhang2021labeling}; and
Neo-GNN, BUDDY, NCN/NCNC, and LPFormer make related structural signals scalable
\citep{yun2021neognn,chamberlain2023buddy,wang2024ncn,shomer2024lpformer}.
These methods establish that topology is useful, but they read it from a graph
that already contains the query endpoints.

\paragraph{Inductive and cold-start link prediction.}
GraphSAGE introduced transferable neighborhood aggregation
\citep{hamilton2017graphsage}, while Graph2Gauss maps attributes to inductive
embeddings \citep{bojchevski2018graph2gauss}. Graph-free students can distill a
graph model into an attribute-only predictor
\citep{zhang2022glnn,zheng2022coldbrew}, and cold-start methods can attach a new
node to an observed graph. These settings differ from the mutually unseen case:
our endpoints have no observed incident edges and no inference graph is supplied.
Direct endpoint-feature scorers satisfy that restriction, but generally score
each pair independently.

\paragraph{Training-time topology transfer.}
Attribute-to-embedding mappings and teacher--student objectives use a training
graph as supervision while keeping inference graph-free. This family is the
closest protocol precedent because structural information can be absorbed into
parameters without exposing test topology. Our question is narrower: whether a
query-conditioned topology representation improves both the binary decision and
the graph assembled from all such decisions.

\paragraph{Graph structure learning and latent inference.}
Graph structure learning estimates an adjacency for a downstream task. LDS uses
bilevel optimization over edge variables \citep{franceschi2019lds}; IDGL
alternates metric learning and graph refinement \citep{chen2020idgl}; and
NodeFormer and DGM provide scalable differentiable construction
\citep{wu2022nodeformer,kazi2023dgm}. Most work optimizes one closed node set or
refines an observed graph. It does not infer query-local context from two unseen
endpoint features for an unchanged binary edge target.

\paragraph{Generated structure as intermediate context.}
Autoregressive and diffusion graph models provide machinery for constructing
sparse adjacency \citep{you2018graphrnn,vignac2023digress}. Neural relational
inference demonstrates the broader pattern of latent structure serving a
downstream prediction \citep{kipf2018nri}. In our setting, generation is only one
candidate representation family. Any inferred structure must remain upstream of
the edge classifier and must not introduce node identities or target-graph
observations unavailable under the endpoint-only contract.

\paragraph{Assembled-graph evaluation.}
Graph-generation evaluation commonly compares degree, clustering, and spectral
descriptors. Prior work shows that such statistics have distinct sensitivities
and should not be collapsed into one score
\citep{thompson2022evaluation,obray2022evaluation}. We accordingly report GS,
RD, and three MMD ratios separately, together with edge-level discrimination and
calibration metrics.

\paragraph{Positioning.}
Table~\ref{tab:positioning} separates methods by inference-time topology source
and output contract. The target setting is edge prediction for two unseen nodes,
conditioned on topology inferred without access to the target graph.

\begin{table}[t]
\caption{Positioning of related families.}
\label{tab:positioning}
\begin{center}
\small
\begin{tabular}{lccccc}
\toprule
\textbf{Family} & \textbf{Unseen} & \textbf{Needs target} & \textbf{Topology-} & \textbf{Edge} & \textbf{Graph} \\
 & \textbf{nodes} & \textbf{graph} & \textbf{conditioned} & \textbf{output} & \textbf{metrics} \\
\midrule
Endpoint-feature scorers             & yes     & no      & no      & yes & no \\
Observed-topology predictors         & partial & yes     & yes     & yes & no \\
Training-time topology transfer      & yes     & no      & partial & yes & no \\
Graph structure learning             & varies  & usually & yes     & no  & no \\
Graph generation                     & ---     & no      & ---     & no  & yes \\
\midrule
\textbf{This work}                   & \textbf{yes} & \textbf{no} & \textbf{yes} & \textbf{yes} & \textbf{yes} \\
\bottomrule
\end{tabular}
\end{center}
\end{table}
